
\AddSymbolsC
\pagecolor{purple!10!gray}
\color{white}
\quad\quad \lettrine[,lraise=0.2,lhang=0.5]{\textbf{铃}}{木}和子躺倒在了树下。她感到头痛不已，用手指用力地推着自己的太阳穴。

猫站在一旁，不停地围着树奔走；走毕，它便回到铃木和子身边，把猫掌轻轻搭在了她的额头上。

而铃木和子则把猫掌温柔地拨开；她只想安静地休息。她侧过身来，看着伊藤佐一留下的脚印，更觉头痛欲裂了。

她忘了上次头痛她是怎么缓解的了——或许她是昏睡过去而不知情的；而这次可能也一样。

不知不觉，傍晚的脚步声响了。街灯像往常一样，同时亮起。铃木和子的的确确昏睡过去了。直到太阳彻底落山，铃木和子才苏醒过来。

\twkkai{到底是怎么回事呢？}她坐了起来，捂着头，想着一些东西。

她想起了自己前几天缺课前在教授课上记下的笔记；可是她总感觉笔记里缺了点什么。她突然有一种奇怪的感觉，因为她似乎察觉到了某些事情——或许和伊藤佐一有关——的异常；但她说不出来。

她站起了身，拍了拍衣服，清理掉身上的粉紫色花瓣。猫在一旁盘坐着，依旧静静地盯着铃木和子。

铃木和子悄然回头，看着猫的圆形的瞳孔——她猛地瞧见里面站着一个人。

猫呜咽着；铃木和子则蹲下，伸出手摸了摸猫，轻抚着猫的身体。

\twkkai{谢谢你关心我——可惜我要走了……}

铃木和子用柔和的语气说着；她的手仿佛离不开猫了，只是机械地继续轻抚着。但她知道，是时候离开了。

她在心里与猫默默告别，便再次站起了身，朝地铁站的方向走去。

猫并没有动。

……

地铁站内的人寥寥无几。铃木和子缓慢地走着；之所以选择了慢步行走，是因为她开始觉得周围的世界变得真实起来。她此刻只想要感受——而前两次头痛成功地将她带回了现实，让她得以看清平时回家时不会留意到的细节。

不过，今天，她并不打算太早回家。

一切都过于反常了；猫，奖章，头痛——没有人能向她解释这一切的怪异。她希望她不是世界上唯一一个觉察到怪异的人。而她把最大的希望寄托在了伊藤佐一身上。

她想着，又转而看着地铁玻璃门里的自己。那个自己，被飞驰而过的列车穿透、打破，变得支离破碎；转而变成的，是众生的脸。那是一张张受了奔波的伤害而痛苦的脸，它们低垂着，仿佛随时要融化。

地铁门开了；人们纷纷走出车厢。厢内瞬间变得空荡荡的，十分清冷，还带着雨水的味道。铃木和子知道，就在刚才等车的十几分钟之间，下了一场雨。

她突然担心起那只猫。猫的毛发，猫的爪子，猫的温度——都是那么真实。但她也时而疑惑——猫是否真的比周遭的世界更为真实。

她走进了车厢，找了最近的一个空座位，坐了下来。车厢门履行着它的职责，缓缓地关上了。车厢摇晃着前进；窗外闪过一幅幅炫目的广告牌，刊登着不同的产品、人物。铃木和子仿佛没有见过这些，但她此时也并不感兴趣。

她能逐渐感受到一股向上升的加速度。列车离开了地底隧道，来到了城市上方的轨道桥上。摩天楼拔地而起；霓虹灯的光芒在空气中散射着；汽车呼啸而过。

铃木和子越来越怀疑这一切的真实性了。\twkkai{美妙得可怕；}她想。她的双手撑着冰冷的不锈钢座椅；双脚自然地悬挂着，没有触碰地面。

没有广告牌；没有了人们。只有冰冷的建筑。

这让她想起了猫。猫的毛发，猫的爪子，猫的温度——对，尤其是温度——那是让她唯一感到安心的。

可是，每次猫的到来，都会不可避免地伴随着她的头痛——尽管只有两次，但已经足以让她警惕。

她仔细地想着。她把身子挪前，好让双脚触地，不再跟着列车的摆动而晃荡——仿佛这样能让她的内心有着落。

而后，她才意识到一些东西——猫总是衔着一块什么。她也并没有留意；她恨自己没有去留意而轻易地让伊藤佐一拿回——仿佛那是属于她的东西。她开始不那么怀疑猫了，又转而相信起猫的美好意义。

车厢中，依旧是那么地寂静。只有列车的呜呜声。或许还有铃木和子的呼吸声。

而铃木和子在一呼一吸间，决定要做点什么。

\begin{center}
——————
\end{center}

夜晚的图书馆里，人影寥寥。只有清脆的脚步声回荡在图书馆的角落。铃木和子看着那几个人的背影，目光逐渐转移到他们的脚跟。那是声音的来源。她试图启发着自己。

可是她始终不知道，自己头痛的来源在哪里。

不过，现在她已经不头痛了；她便趁着这个间隙，拿出了许久不见的蓝牙耳机。

音乐在耳畔奏起；铃木和子的脚步随着音乐的推进而时缓时快。

而慢步走时，她又想起那个人。

那个站在猫眼睛里的人——他的形象是多么地清晰，可是她无论如何也无法说出关于这个人的任何信息。

\twkkai{我见过吗？}

楼梯逐渐逼近；铃木和子抬起头，好让自己看得清阶梯。她到了二楼平台。

她突然停了下来。不远处，正站着一个身影，在书架前徘徊着。他低垂着头，手里拿着一张纸，仿佛在斟酌什么。

铃木和子不知道，她见到的那个身影，和刚才她想到的身影，是否有联系。她摘下了耳机，无所适从。她此刻突然是那么希望自己头痛，从而能躺倒在地上，睡醒后一切皆正常。

那个身影依旧时远时近。铃木和子犹豫着，缓缓地走向了那个身影对面的书架。

书架旁的空间本来昏暗；她一走进那片漆黑，头顶上的灯便突然亮起。铃木和子瞬间听到一阵沙沙声。她决定躲起来。

可是灯光下，何处可躲呢？那个身影早已走近，逐渐显现出立体的轮廓。

是伊藤佐一。他率先出声。

\twkkai{又见面了……铃木同学？}他心里是那么确信，可是他故意装得糊涂。

铃木和子突然不敢吭声；就在那一刻，她仿佛又头痛了起来——虽然实际上并没有。她想到了什么——仿佛她心底里有什么秘密不可以告诉他，又或者他将要告诉她一些可怕的真相；她更怕的是后者。

她的身子挪后了一步。伊藤佐一见此，便换了语气。

\twkkai{你……也是过来复习的吗？}

他知道，这个话题似乎轻松点；可铃木和子不那么觉得。她依旧躲得远远的；她害怕。

伊藤佐一开始疑惑起来。他的直觉让他不可抗拒地看出了她眼中的抗拒。他这才意识到自己的尴尬。

\twkkai{好吧，打扰了……}

伊藤佐一转过身去。他没有回头。他走回了原来的位置，继续拿出那张纸。\twkkai{明明傍晚的时候不是这样的……；}他想着。

铃木和子心里满是愧疚。她知道，就在刚刚，她把自己的头痛全都归因到了眼前站着的这个人身上。可她也打心底里不愿靠近他——因为她突然才意识到，刚才在地铁上做出的判断，是彻头彻尾的错误。他刚才轻佻的态度表明了他觉察不到这个世界的怪异；相反，他似乎是导致她周遭世界变得怪异的原因。

她越来越相信这一点，以至于她又走到了更远的一个书架。她也看起了书。

可是没过一会儿，她便听到清脆的“咣当”一声。她立即探出了头。

是那枚奖章。它滚落在远处的椅子一角，而伊藤佐一慌张地跑去捡起；待他捡起后，他便看到铃木和子远处的目光。铃木和子把头缩了回去；而伊藤佐一的脸又“唰”地一下红了。他急匆匆地回到原位，又拿出了那张纸。

月亮的光打在伊藤佐一的纸上，也打在铃木和子的书上。

……

已经开始逼近八点了。伊藤佐一看了看自己的手表。

\twkkai{该回家了；}他想。他折好手中的纸，离开了书架。他走下了楼梯。

\twkkai{等等！}

伊藤佐一回过了头。只见铃木和子合上了书，站了起来，拍了拍自己的裙子。

\twkkai{我有话对你说；}她看着伊藤佐一的双眼，并不回避。

平时连女生眼睛都不敢对视的伊藤佐一，这次冒着被铃木和子误会的风险，也选择盯着铃木和子。

\twkkai{请说；}他低声说着，而后吞了吞口水。

\twkkai{我请假的那天，教授是不是说要考试？}铃木和子问道，目光没有移走。

伊藤佐一没有想到她竟然只问这样的问题。

\twkkai{是的……}

\twkkai{那我能借你的笔记抄一抄吗？}

伊藤佐一呆住了一霎那。他想起教授的告诫。

\twkkai{……切忌在课堂上拿出来并提问！}

他的双手开始变得无处安放；他不敢把手插进口袋里，怕引起铃木和子的注意。

\twkkai{嗯……等明天吧？今天已经晚了；}伊藤佐一把自己的支吾化成了思考的痕迹。

他看着脚下。他发现书架的影子较刚才已经缩短了一点了。可铃木和子的目光自始至终都没有变化——反倒更坚定了。

伊藤佐一看着铃木和子的双眸——他猛地瞧见里面站着一个人。

而铃木和子不知道自己发起的对话应该如何收尾。但她竟说了那个词。

\twkkai{不行。}

她没想到自己竟略带委屈地说出了这个词。伊藤佐一更是意料不到她会说出这个词——他们之间并没有那么熟络；她没有理由如此要求他。

但影子越缩越短，书架底部的那块砖被月亮照得越来越亮。

\twkkai{走了！不要站在这里了！}

楼梯间里出现了一位保安。

\twkkai{非常抱歉，我……我还是明天给你吧……}

伊藤佐一扭头看向保安，又看向铃木和子，说着。

保安向他招手示意，他见此，便开始往楼梯口逃窜，留下铃木和子一人。

……

\begin{center}
——————
\end{center}

天气依旧。依然晴朗。

伊藤佐一照常迈进课室。周围的同学推搡着他，而他混在人群中，经过了讲台旁的教授。教授没有看他。

伊藤佐一坐回了那个位置——那个他曾经感到不安的位置。坐下后，他立即回头看了一眼。

铃木和子正看着他。她的眼神里，没有愠怒，没有责备，只有平静。

可是这平静让伊藤佐一感到极度不安。他扭回了头，恨自己昨晚当着她的面逃走。

可是他又想起，他们之间并不熟络——甚至可以说，在此之前并不相识；这样来看，他并没有做错什么。

他全身发麻，但感觉比刚才好一点了。他肯定着自己的逻辑。

课室里坐满了人。讲台上的教授整理好讲义，抬起了头。

\twkkai{好了，同学们——今天考试！}

周围一片寂静。仿佛那些学生都知道今天考试。

而伊藤佐一才想起，的的确确是今天考试。他的冷汗又冒出来了。他动都不敢动；但他手里没有笔，必须要从包里拿。

考卷从下面的座位一层一层递上来。伊藤佐一只好迅速地回头拿起书包，拉开拉链，拿出笔。就在拿起书包的一瞬，他瞟了一眼。

自然是瞟向铃木和子的。

她正趴在桌子上。她在哭。

\twkkai{咳！拿试卷了！}

一叠试卷“哗啦”一声掉在了伊藤佐一面前。伊藤佐一此时顾不得那么多了，轻轻拿了一张，便又往后传。

开考了。

伊藤佐一看向试卷。铃木和子的抽泣无法逃离他的脑海。

\twkkai{可是又能怎样呢……等等？}

他猛地瞧见试卷上的证明题。那正是奖章上刻的铭文：
\begin{gather*}
L(s,f\otimes g)=\zeta(2s-2k+2)\sum_{n=1}^{\infty}\frac{a(n)b(n)}{n^s},\\
\text{\twkkai{其中} }f(z)=\sum_{n=1}^{\infty}a(n)q^n,\ g(z)=\sum_{n=1}^{\infty}b(n)q^n,\ q=e^{2\pi iz},\ \Im(z)>0.
\end{gather*}\index{ljuanji@$L(s,f\otimes g)$（Rankin-Selberg $L$-函数）}

伊藤佐一依稀记得教授的讲解。但是铃木和子伏在桌面上的身影依然萦绕在他的脑海中，干扰着他。

他不知道该考试好，还是该想想考完后怎么安慰她好。

他看见教授走出了课室。一个念头突然出现。

……

考试结束了。伊藤佐一交了卷后，才逐渐从数学世界的沉浸中抽离出来。

他放好了笔。回头一看，铃木和子的身影不见了。他又把头扭向另外一边；铃木和子刚好经过了他，正走下他旁边的课室阶梯。

她没有理他。

伊藤佐一看着铃木和子的背影，更加笃定地认为，自己刚刚做的决定是正确的。

\begin{center}
——————
\end{center}

\twkkai{好了，同学们。}

教授用手指敲了敲讲台。显然，过了仅仅一天，他便改完了试卷。

\twkkai{大家考得都挺不错。除了两位同学。}

\twkkai{伊藤佐一和铃木和子，你们两个等下到我的办公室……}

课室内传来一阵轻微的讨论声。而伊藤佐一看着众人的反应，知道自己的行动奏效了。但他还需要铃木和子的反应。

……

教授的课对伊藤佐一来说变得越来越无聊。伊藤佐一看着手表上的指针，发着呆。

终于等到下课了。依旧是人群涌动。伊藤佐一知道什么将要发生；但他不敢确定。

他害怕。

他害怕铃木和子也像他那样逃走——那样他就会永远欠她一个解释。

可是就在这一刻，他看到一个熟悉的背影走在他前面。正是铃木和子。她将去往何方？

伊藤佐一内心有跟上她的冲动，但仔细地斟酌过后，他决定还是缓缓地走好自己脚下的路。他走下了一级又一级阶梯，走出了课室门口。

走廊里没有了铃木和子的身影。伊藤佐一开始担忧了。

他来到了办公室。教授正坐在椅子上，双手交叉，搭在大腿上，闭眼休息着。他没有听到伊藤佐一的脚步声。显然，教授已经熟睡。

这是伊藤佐一有史以来过得最煎熬的时光。他就那么站着，可他内心是多么想劝他去周围走走。他转而看向教授身后的黑板。那里正写着一行史诗般让人惊心动魄的推导。

他走近了，看着。
\begin{gather*}
\sum_{n=1}^{\infty}a(n)b(n)n^{-s}=\prod_p\left(\sum_{l=0}^{\infty}a(p^l)b(p^l)p^{-ls}\right).
\end{gather*}

他正看得入迷，却没看到铃木和子早已站在门口；他这才抬起头，偷瞄了她一眼。一瞬间，难熬的时光似乎终于过去了；黑板上的数学推导变得像水里的鲤鱼一般鲜活起来。

伊藤佐一这才意识到——自己喜欢上了铃木和子。而铃木和子此时走得更近了；她把身子往桌子上一弯，放好了一叠试卷。

\twkkai{原来教授安排她整理试卷了；}伊藤佐一松了一口气。

\twkkai{教授；}铃木和子开口了，似乎毫不客气——仿佛是来和教授谈判的。

教授睁开眼睛，坐起身子。他翻找着桌面上的试卷，抽出两张——正是伊藤佐一和铃木和子的。

\twkkai{自己解释一下吧——伊藤同学，你先来；}教授平静地说着。

伊藤佐一看向铃木和子。她没有看他，而是盯着黑板上那条公式发呆——正是那条仿佛变成了鲤鱼的公式。

\twkkai{我……教授，我当时没有搞懂题目……}

教授盯紧了伊藤佐一。但他那锁紧的目光中似乎包含着一种对伊藤佐一的提醒。他没有说话，但也没有让沉默维持太久。

\twkkai{那铃木同学呢？}

\twkkai{我忘了复习，然后……最近头痛还没有恢复；}铃木和子把目光从公式上移开，看着自己的手，腼腆地说着。她没有了刚才的灵气。

伊藤佐一第一次打心底里心疼铃木和子。他开始猜测她头痛的原因。他又想起教授的那双目光。

\twkkai{好，没事；我和其他教授不一样，不会打着某种名义批评你们——我从不会强迫什么事情；但是，伊藤同学……}

教授左右旋转着座椅说，又转向伊藤佐一。\twkkai{……我给你一个小任务：每天辅导一下铃木同学——十分钟也好啊；}教授再次盯紧了伊藤佐一。

而铃木和子一直低垂着头，全程都没有留意到伊藤佐一和教授的舞台剧。

伊藤佐一不断地受着教授的提醒，仿佛终于想起了什么。

\twkkai{是奖章！是吗？}他想着。

\twkkai{好的；}他答着。

\twkkai{那你们可以离开了；}教授轻轻挥了挥手，眼里带着某种亲切。

而伊藤佐一早就在刚才的短暂对话中下定了某种决心——他把整块黑板的推导都大致记了下来。

